\chapter*{DANH MỤC TỪ VIẾT TẮT VÀ KÝ HIỆU}
\addcontentsline{toc}{chapter}{Danh mục từ viết tắt và ký hiệu}

\begin{longtable}{@{}p{3.2cm}p{11.3cm}@{}}
\toprule
\textbf{Viết tắt} & \textbf{Ý nghĩa} \\
\midrule
\endhead
MOGVRPTW & Multi-Objective Green Vehicle Routing Problem with Time Windows: bài toán định tuyến xe xanh
           đa mục tiêu có khung thời gian \\
NSGA-II  & Non-dominated Sorting Genetic Algorithm II: giải thuật di truyền sắp xếp không trội, phiên bản 2 \\
NN       & Nearest Neighbor: láng giềng gần nhất \\
TW       & Time Window: khung thời gian; trong báo cáo dùng để chỉ cách khởi tạo theo thứ tự khung giờ \\
LS       & Local Search: tìm kiếm cục bộ \\
OX       & Order Crossover: lai ghép thứ tự \\
DP       & Dynamic Programming: quy hoạch động (cách cắt tuyến tối ưu của TV4) \\
HV       & Hypervolume: thể tích vùng bị trội \\
IGD      & Inverted Generational Distance: khoảng cách thế hệ đảo \\
SP       & Spacing: độ đều của tập nghiệm \\
SC       & Set Coverage: độ phủ tập \\
RNG      & Random Number Generator: bộ sinh số ngẫu nhiên \\
TV1--TV4 & Thành viên 1 đến thành viên 4 của nhóm \\
\midrule
\textbf{Ký hiệu} & \textbf{Ý nghĩa} \\
\midrule
$n$                     & số khách hàng (không tính kho, kho có chỉ số 0) \\
$\pi$                   & hoán vị khách hàng, là nhiễm sắc thể của một cá thể \\
$Q$                     & tải trọng tối đa của một xe (kg) \\
$K$                     & số xe tối đa của đội xe (\code{fleet\_size}) \\
$d_{ij}$                & khoảng cách Euclid giữa điểm $i$ và $j$ (km) \\
$[ET_i, LT_i]$          & khung thời gian chấp nhận được của khách $i$ \\
$[ETa_i, LTb_i]$        & khung thời gian kỳ vọng của khách $i$, nằm bên trong $[ET_i, LT_i]$ \\
$f_1, f_2, f_3$         & chi phí (CNY), khí thải CO\textsubscript{2} (kg), mức hài lòng trung bình (0--1) \\
$\mathbf{g}$            & vectơ mục tiêu cực tiểu hoá $(f_1, f_2, 1-f_3)$ \\
$N_{pop}$, $G$          & kích thước quần thể, số thế hệ \\
$p_c$, $p_m$, $p_{ls}$  & xác suất lai ghép, đột biến, chạy tìm kiếm cục bộ \\
$\mathbf{w}$, $\mathbf{s}$ & vectơ trọng số và vectơ hệ số chuẩn hoá trong tiêu chí tổng có trọng số \\
\bottomrule
\end{longtable}
